\section{Interpretation, limitations, and next steps}
The view-line idea survives the addition of a consistent local perturbation.
At a fixed reception event, different directions connect the observer to
different spatial locations and earlier emission events. On the expanding
branch, those events can be drawn on progressively smaller background
layers. This is the precise meaning of ``looking inward through cosmic
history.'' It identifies a direction in the diagram's time coordinate,
not a preferred spatial direction for a telescope.

The local-mass calculation also sharpens the original picture. The
spatial deformation is controlled by $1-2\ph$, clock rates by $1+2\ph$,
and null propagation by their ratio. Their relationship comes from the
metric and field equations. It cannot be determined from the apparent
depth of an arbitrarily drawn depression. In particular, the first-order
spatial embedding has an outward core deformation and a reduced clock
rate simultaneously. A clock that accumulates less proper time remains
on a future-directed worldline; it has not moved to an earlier cosmic
slice.

The model's physical content is deliberately restricted. It is a linear
growing dust perturbation, with a specified compensation profile on a
closed background. It has no horizon and no central singularity.
Extrapolating its lapse toward zero would leave the approximation's
domain of validity. A later black-hole extension must use a suitable
strong-field metric, matter model, matching conditions, and a regular
choice of coordinates. Non-flat McVittie constructions illustrate that
embedding a mass in a closed expanding cosmology is a separate problem
with substantial global constraints \citep{Nolan2017}.

There are further representational limits. The source matching, spatial
embedding, and separated-clock comparisons are made in an explicitly
chosen gauge. Observable angles or frequency ratios require physical
observer frames; coordinate radius or timing alone is not an invariant.
Our matched-source angle uses such a frame, while the clock-offset and
conformal-delay plots retain their stated background-relative meaning.
The code's geodesics include some nonlinear dependence on a metric whose
Einstein dynamics are only linear. Amplitude scaling exposes this
distinction but does not supply the missing higher-order dynamics.

The strongest immediate use is therefore a reproducible comparison:
the same mass perturbation gives a spatial metric, a lapse, and view
lines, with their different graphical meanings kept explicit. No claim
of new cosmological dynamics or observational evidence follows. The
accompanying \emph{Comparing Cosmological Diagrams} note evaluates the
same events in four maps and identifies compression and time-folding
limits. It finds no additional physical information in the nested map;
an explanatory benefit from pairing representations remains an untested
reader-level hypothesis.

Useful next steps are to evolve emitting and receiving dust tracers,
include their Doppler terms, calculate a full angular-distance/Jacobi
map, vary curvature and background composition, and compare the weak-field
result with an exact mass-in-cosmology construction. A pedagogical study
could separately test whether readers better distinguish spatial
separation, elapsed proper time, and past-light-cone observation with
these paired diagrams.
